%Chargement des paquets
\usepackage{amsmath}
\usepackage{amsfonts}
\usepackage{amssymb}
\usepackage{enumerate}
\usepackage{mathtools}
\usepackage{bbm}
\usepackage{xparse, etoolbox}
\usepackage{enumerate}
\usepackage{enumitem}
\usepackage{mathabx}
\usepackage{babel}[french]

%environment
\usepackage{amsthm}
\usepackage{keytheorems}
\usepackage{hyperref} 

%Interface théorème
\newkeytheorem{lem}[
	name = Lemme
]

\newkeytheorem{prop}[
	name = Proposition
]

\newkeytheorem{thm}[
	name = Théorème
]

\newkeytheorem{coro}[
	name = Corollaire
]

\renewcommand*{\proofname}{Démonstration}

\theoremstyle{definition}
\newtheorem{defn}{Définition}

\theoremstyle{definition}
\newtheorem*{nota}{Notation}

\theoremstyle{definition}
\newtheorem*{limi}{Liminaire}

\theoremstyle{remark}
\newtheorem{rmq}{Remarque}

\theoremstyle{plain}
\newtheorem*{fait}{Fait}


%Ensembles de nombres
\newcommand{\N} {\mathbb{N}}
\newcommand{\Ne}{\N^\ast}
\newcommand{\Z} {\mathbb{Z}}
\newcommand{\Q} {\mathbb{Q}}
\newcommand{\R} {\mathbb{R}}
\newcommand{\Rb}{\overline{\mathbb{R}}}
\newcommand{\Rp}{\R_+}
\newcommand{\Rm}{\R_-}
\newcommand{\K} {\mathbb{K}}
\newcommand{\Ket} {\mathbb{K^{\ast}}}
\newcommand{\Cx}{\mathbb{C}}

%Ensembles, tribus
\newcommand{\A}{\mathbb{A}}
\renewcommand{\P}{\mathcal{P}}
\newcommand{\T}{\mathcal{T}}
\newcommand{\F}{\mathcal{F}}
\newcommand{\C} {\mathcal{C}}
\newcommand{\ouv}{\mathcal{O}}
\newcommand{\B}{\mathcal{B}}
\newcommand{\M}{\mathcal{M}}
\newcommand{\etag}{\mathcal{E}}
\newcommand{\etagOp}{\etag^{+}(\Omega)}
\newcommand{\etagO}{\etag(\Omega)}
%%Lp avant quotientage
\newcommand{\Lic}{\mathcal{L}^1}
\newcommand{\Lpc}{\mathcal{L}^p}
\newcommand{\Linfc}{\mathcal{L}^{\infty}}
\newcommand{\Lifull}{\Lic(\Omega, \B, m)}
\newcommand{\Lpfull}{\Lpc(\Omega, \B, m)}
\newcommand{\Linffull}{\Linfc(\Omega, \B, m)}
%%Lp après quotientage
\newcommand{\Li}{L^1}
\newcommand{\Ld}{L^2}
\newcommand{\Lp}{L^p}
\newcommand{\Linf}{L^{\infty}}
\newcommand{\Listd}{\Li(\Omega, m)} 
\newcommand{\Ldstd}{\Ld(\Omega, m)} 
\newcommand{\Lpstd}{\Lp(\Omega, m)} 

%Quelques opérateurs
\DeclareMathOperator{\codim}{codim}
\DeclareMathOperator{\rg}{rg}
\DeclareMathOperator{\tr}{tr}
\DeclareMathOperator{\vect}{Vect}
\DeclareMathOperator{\ima}{Im}
\DeclareMathOperator{\id}{Id}
\DeclareMathOperator{\sign}{sign}
\DeclareMathOperator{\ext}{ext}
\newcommand{\ind}{\mathbbm{1}}
\newcommand*{\spr}[2]{\left(\,#1\,\middle|\,#2\,\right)}
\newcommand{\interior}[1]{%
  {\kern0pt#1}^{\mathrm{o}}%
}
\DeclareMathOperator{\card}{card}
\DeclareMathOperator{\ppcm}{ppcm}

%Commandes de pure flemme
\newcommand{\veps}{\varepsilon}
\newcommand{\intab}{\int_{a}^{b}}
\newcommand{\intR}{\int_{-\infty}^{\infty}}
\newcommand{\intinf}{\int_{- \infty}^{+ \infty}}
\newcommand{\equi}{\Leftrightarrow}
\newcommand{\Kev}{$\K$-espace vectoriel }
\newcommand{\Kevs}{$\K$-espaces vectoriels }
\newcommand{\Rev}{$\R$-espace vectoriel }
\newcommand{\Revs}{$\R$-espaces vectoriels }
\newcommand{\Cev}{$\Cx$-espace vectoriel }
\newcommand{\Cevs}{$\Cx$-espaces vectoriels }
\newcommand{\evn}{(E, \norm{.})}
\newcommand{\interf}[2]{[#1,#2]}
\newcommand{\limb}{\lim\limits_}
\newcommand{\lebn}{\lambda_n}
\newcommand{\pp}{\text{~presque partout}}
\newcommand{\derivp}[2]{\frac{\partial #1}{\partial #2}}
\newcommand{\famiu}{(u_i)_{i \in I}}
\newcommand{\sev}{\text{~s.e.v.}}
\newcommand{\Mnp}{M_{n,p}(\K)}
\newcommand{\Mn}{M_{n}(\K)}
\newcommand{\tq}{\text{~t.q.~}}
\newcommand{\mq}{~montrer~que~}
\newcommand{\et}{\quad \text{et} \quad}
\newcommand{\ou}{\text{~ou~}}
\newcommand{\Kx}{\K[X]}


%Commandes de gars chelou
\renewcommand{\subset}{\subseteq}

%Commande restriction, ty duckduckgo
\def\restr#1#2{\mathchoice
              {\setbox1\hbox{${\displaystyle #1}_{\scriptstyle #2}$}
              \restraux{#1}{#2}}
              {\setbox1\hbox{${\textstyle #1}_{\scriptstyle #2}$}
              \restraux{#1}{#2}}
              {\setbox1\hbox{${\scriptstyle #1}_{\scriptscriptstyle #2}$}
              \restraux{#1}{#2}}
              {\setbox1\hbox{${\scriptscriptstyle #1}_{\scriptscriptstyle #2}$}
              \restraux{#1}{#2}}}
\def\restraux#1#2{{#1\,\smash{\vrule height .8\ht1 depth .85\dp1}}_{\,#2}} 

%Quelques normes
\DeclarePairedDelimiter\abs{\lvert}{\rvert}
\DeclarePairedDelimiter\labs{\bigl\lvert}{\bigr\rvert}
\DeclarePairedDelimiter\norm{\lVert}{\rVert}
\DeclarePairedDelimiterXPP\normi[1]{}\lVert\rVert{_1}{\ifblank{#1}{\:\cdot\:}{#1}}
\DeclarePairedDelimiterXPP\normd[1]{}\lVert\rVert{_2}{\ifblank{#1}{\:\cdot\:}{#1}}
\DeclarePairedDelimiterXPP\normp[1]{}\lVert\rVert{_p}{\ifblank{#1}{\:\cdot\:}{#1}}
\DeclarePairedDelimiterXPP\normq[1]{}\lVert\rVert{_q}{\ifblank{#1}{\:\cdot\:}{#1}}
\DeclarePairedDelimiterXPP\norminf[1]{}\lVert\rVert{_\infty}{\ifblank{#1}{\:\cdot\:}{#1}}

%Bracket en angle
\DeclarePairedDelimiter\angl{\langle}{\rangle}

%Set, parenthèses
\newcommand{\set}[1]{\{ #1 \}}
\newcommand{\bigset}[1]{\bigl\{ #1 \bigr\}}
\newcommand{\Bigset}[1]{\Bigl\{ #1 \Bigr\}}
\newcommand{\bigpar}[1]{\bigl( #1 \bigr)}
\newcommand{\Bigpar}[1]{\Bigl( #1 \Bigr)}

%Opitmisation des opérateurs de comparaison
\DeclarePairedDelimiter\tio{o (}{)}
\DeclarePairedDelimiter\gro{O (}{)}

\newcommand{\Bfr}{\widebar{B_r}}

\pagestyle{fancy}

\fancyhead[C]{}
\fancyhead[R]{}

\begin{document}

\underline{Source :} 
\begin{itemize}
\item Rouvière - Petit guide - page 223
\item Gourdon - Analyse - page 341
\end{itemize}

\underline{Leçons :}
\begin{itemize}
\item 215 : Applications différentiables définies sur un ouvert de Rn. Exemples et applications.
\item 214 : Théorème d’inversion locale, théorème des fonctions implicites. Illustrations en analyse et en géométrie.
\end{itemize}

\begin{nota}
On note $\normt{\cdot}$ la norme définie sur $\mathcal{L}(\R^n, \R^n)$ par :
\[ \forall u \in \L(\R^n, \R^n), \quad \normt{u} = \sup_{\norm{v}=1} \norm{u(v)} . \]
\end{nota}

\begin{thm}
Soit $U$ un voisinage ouvert de $\R^n, n \geq 1$ et $f:U \to \R^n$ une application de classe $\C^1$.
On suppose qu'il existe $a \in U$ tel que $df_a$ soit un isomorphisme bicontinu de $\R^n$.
Alors il existe un voisinage ouvert $V$ de $a$ et un voisinage ouvert $W$ de $f(a)$, tous deux dans $\R^n$ 
et tels que $f$ soit un $\C^1$-difféomorphisme de $V$ sur $W$.
\end{thm}

\begin{proof}
La démonstration est divisée en plusieurs parties pour simplifier sa lecture.
\begin{enumerate}[label=(\roman*)]

\item
On montre dans un premier temps que le problème peut être ramené au cas $a=f(a)=0$ et $df(0)=\id$. 
En effet, on peut toujours substituer à $f$ la fonction $x \mapsto df(a)^{-1}(f(a+x) - f(a))$,
i.e. on \og centre \fg{} la fonction en $a$ (ce qui est possible car $f$ est défini sur un ouvert contenant $a$) et on étudie la variation
de $f$ autour de $a$ (soit $f(a+x)-f(a)$). De plus, comme $df(a)$ est une bijection on peut appliquer à ces variations la transformation
linéaire inverse $df(a)^{-1}$. Ainsi, au voisinage du point $a$, la fonction nouvellement obtenue est l'application linéaire identique plus
un petit reste négligeable devant la norme de la variation à $a$. \\

\item
Ces hypothèses étant prises, on pose, pour $y \in \R^n$, $F_y$ l'application définie sur $U$ par :
\[ F_y(x) = x + y - f(x) . \]
Soit $\veps \in ]0,1[$, on va montrer qu'il existe $r>0$, indépendant de $y$, tel que :
\[ \normt{dF_y(x)} \leq \veps \text{ pour } \norm{x} \leq r . \]
L'application $F_y$ est de classe $\C^1$ et on a $dF_y(x) = \id - df(x)$. 
On en déduit que $dF_y(0) = 0$ et, par continuité, qu'il existe $r$ tel que $\norm{x} \leq r \Rightarrow \normt{dF_y(x)} \leq \veps$.
De plus, comme $dF_y(x)$ est indépendante de $y$, le choix de $r$ l'est aussi. \\

\item
Notons $B_r$ la boule ouverte $B(0, r)$ de $\R^n$, $\Bfr$ sa fermeture et $W$ l'ensemble des $y \in \R^n$ tels que :
\[ \norm{y} < (1-\veps)r . \]
On va montrer que $F_y(\Bfr) \subset B_r$ pour $y \in W$.
Soit $x \in \Bfr$, on a vu que $\normt{dF_y(x)} \leq \veps$, donc, par l'inégalité de la moyenne :
\[ \norm{F_y(x)} \leq \norm{F_y(0)} + \norm{F_y(x) - F_y(0)} \leq \norm{y} + \norm{F_y(x) - y} < (1-\veps)r + \veps r < r , \]
et ainsi, on a bien $F_y(x) \in B_r$.

\item 
Posons $V = B_r \cap f^{-1}(W)$, on va démontrer que $f$ est une bijection de $V$ sur $W$.
Le point (i) affirme que $F_y$ est $\veps$-contractante sur $\Bfr$ et (ii) affirme que $F_y$ laisse $\Bfr$ stable si $y \in W$.
D'après le théorème du point fixe de Picard, $F_y$ admet un unique point fixe et que ce point fixe est limite de tout orbite suivant $F$.
Par fermeture, ce point fixe est donc un élément $\Bfr$ (obtenu en suivant par exemple l'orbite selon $F$ débuté en 0), 
comme de plus $F_y(\Bfr) \subset B_r$ selon (ii), on déduit que $x \in B_r$.
En résumé, on a démontré qu'il existe un voisinage ouvert de 0, $W$, tel que, tout point $y \in W$ admette un antécédent unique par $f$
dans $B_r$. 
Cependant, certains points de $B_r$ pourraient avoir leurs images par $f$ en dehors de $W$, on pose donc $V = B_r \cap f^{-1}(W)$,
$V$ contient bien 0 ($f(0)=0$) et est ouvert. On a $f(V) \subset W$ et tout point de $W$ admet un antécédent unique dans $V$, 
ce qui démontre que $f$ est une bijection de $V$ dans $W$.

\item
On veut maintenant montrer que l'image par $f$ de $\Bfr$ est telle que :
\[ (1-\veps)\Bfr \subset f(\Bfr) \subset (1+\veps)\Bfr . \]
Soit $x \in \Bfr$, on remarque d'abord que, pour $y$ fixé on a :
\begin{align*}
\norm{f(x)} & = \norm{F_y(x) - F_y(0) - x} \\
& \leq \norm{F_y(x) - F_y(0)} + \norm{x} \\
& \leq \veps \norm{x} + \norm{x} = (1+\veps) r ,
\end{align*}
selon l'inégalité de la moyenne, ce qui prouve l'inclusion de droite. \\
Par hypothèse $W = B(0, (1-\veps)r) = (1-\veps)B_r$. Or, selon le point (iii), tout $y \in W$ admet un unique antécédent dans $\Bfr$
ainsi, on a nécessairement $W \subset f(\Bfr)$. Comme $f$ est continue et $\Bfr$ compact, on déduit que $f(\Bfr)$ est compact
et ainsi, toute suite convergente de points de $W$ converge dans la fermeture de $W$ et par suite dans $f(\Bfr)$.
Finalement, $(1-\veps)\Bfr \subset f(\Bfr)$, ce qui démontre l'inclusion de gauche.

\item
Montrons maintenant que l'application $f^{-1}$ définie sur $W$ est continue. Soit $y, y' \in W$, il existe $x, x' \in V$ tels que
$f(x)=y$ et $f(x')=y'$. Par construction, $x$ et $x'$ sont les points fixes respectifs de $F_y$ et $F_{y'}$ donc :
\begin{align*}
x-x' & = F_y(x) - F_{y'}(x') \\
& = (F_y(x) - F_y(x')) + (F_y(x')-F_{y'}(x')) \\
& = (F_y(x) - F_y(x')) + (y-y')
\end{align*}
Donc, en utilisant à nouveau l'inégalité de la moyenne :
\[ \norm{x-x'} \leq \veps \norm{x-x'} + \norm{y-y'} , \]
soit,
\[ \norm{f^{-1}(y) - f^{-1}(y')} \leq \dfrac1{1-\veps} \norm{y-y'} , \]
ce qui montre que $f^{-1}$ est lipschitzienne donc continue sur $W$.

\item 
Prenons $x \in \Bfr$ et démontrons que $df(x)$ est inversible. 
Selon le point (ii), on a $\normt{dF_y(x)} = \normt{\id - df(x)} = \veps < 1$.
Par ailleurs, il est clair que $\normt{(\id-df(x))^n} \leq \normt{\id - df(x)}^n < 1$ pour tout $n \in \N$, ainsi la série $\sum \id - df(x)$ converge
absolument. De plus :
\[ \id - (\id - df(x)) \sum_{n=0}^{\infty} (\id - df(x))^n =  \sum_{n=0}^{\infty} (\id - df(x))^n -  \sum_{n=1}^{\infty} (\id - df(x))^n = \id , \]
ce qui montre que $df(x)$ est inversible d'inverse la limite de la série $\sum \id - df(x)$.

\item
On va enfin démontrer que $f^{-1}$ est de classe $\C^1$. Pour cela, on prend $y, y' \in W$ et $x, x' \in V$ tels que $f(x)=y$ et $f(x')=y'$.
On a :
\begin{equation} y - y' = f(x) - f(x') = df(x)(x'-x) + R(x'-x) , \label{eq:1} \end{equation}
avec $R$ une fonction négligeable devant $\norm{x'-x}$ pour $x'-x \to 0$. Donc, pour tout $\eta >0$, on peut trouver $\alpha >0$ tel que :
\[ \norm{x'-x} \leq \alpha \Rightarrow \norm{R(x'-x)} \leq \eta \norm{x'-x} . \]
D'après le point (vi), on a $\norm{x'-x} \leq \dfrac1{1-\veps} \norm{y'-y}$, ainsi en prenant $y, y'$ tels que $\dfrac1{1-\veps} \norm{y'-y} \leq \alpha$
on a :
\[ \norm{R(x'-x)} \leq \eta \norm{x'-y} \leq \dfrac{\eta}{1-\veps} \norm{y'-y}, \]
autrement dit $R$ est négligeable devant $\norm{y'-y}$. On peut donc réecrire (\ref{eq:1}) en :
\[ df(x)(x'-x) = y'-y + o(\norm{y'-y}) , \]
et $df(x)$ étant inversible selon le point (vii), on déduit que :
\[ x'-x = f^{-1}(y') - f^{-1}(y) = df(x)^{-1}(y'-y) + o(\norm{y'-y}) , \]
ce qui établit que la fonction $ f^{-1}$ est différentiable sur $W$ avec, par unicité de la différentielle :
\[ df^{-1}(y) = df(f^{-1}(y))^{-1} . \]
Pour finir, $f^{-1}$ est bien de classe $\C^1$ car son application différentielle est la composée de trois applications continues :
\[ y \mapsto x=f^{-1}(y) \et x \mapsto df(x) \et u \in \mathcal{L}(\R^n, \R^n)^{\ast} \mapsto u^{-1} . \]
\end{enumerate}
\end{proof}

\end{document}